% Story book body — GENERATED, do not edit by hand.
%
% Regenerate with:  python3 tools/make-story.py
% Source of truth:  book/story/pieces.json

\gpart{XVIII}{Komisches}

\gsec{Der Aufkleber}
\levelbadge{A2}

Seit drei Wochen klebt ein Aufkleber auf meinem Küchentisch. Er ist klein und blau. Auf dem Aufkleber steht der Name einer Bananenfirma. Ich habe die Banane gegessen, aber der Aufkleber ist geblieben.

Zuerst habe ich es mit warmem Wasser probiert. Das hat nicht funktioniert. Dann habe ich ein Messer genommen. Meine Frau sagt, dass man einen Tisch nicht mit einem Messer sauber macht. Am Samstag habe ich Öl benutzt, am Sonntag den Föhn. Der Aufkleber ist immer noch da. Er ist jetzt grau, weil ich ihn so oft angefasst habe.

Gestern hat mein Bruder mich besucht. Er hat den Aufkleber in zwei Sekunden entfernt, mit dem Fingernagel. Er hat nichts gesagt. Ich habe auch nichts gesagt. Der Tisch ist an dieser Stelle jetzt ein bisschen heller als überall sonst. Deshalb steht dort eine Schale. Ich kaufe keine Bananen mehr.

\vfill
\begin{glosstbl}
  \textbf{der Aufkleber} & sticker, label & \textbf{kleben} & to stick, to be stuck \\
  \textbf{probieren} & to try & \textbf{benutzen} & to use \\
  \textbf{der Föhn} & hairdryer & \textbf{anfassen} & to touch \\
  \textbf{entfernen} & to remove & \textbf{der Fingernagel} & fingernail \\
  \textbf{die Schale} & bowl & \textbf{die Stelle} & spot, place \\
\end{glosstbl}

\gsec{Der Nachbar auf der Treppe}
\levelbadge{B1}

Am Dienstag hörte ich Herrn Vogler zwei Stockwerke unter mir. Ich erkannte ihn sofort an den Schritten. Damit begann wieder das Problem, das ich seit vier Jahren nicht lösen kann: Wann genau grüßt man einen Nachbarn auf der Treppe?

Grüßt man zu früh, muss man den Rest des Weges schweigend nebeneinander gehen, und das sind elf Stufen. Grüßt man zu spät, wirkt es, als hätte man den anderen absichtlich übersehen. Ich blieb also auf dem Absatz stehen und suchte in meiner Tasche nach dem Schlüssel, den ich in der Hand hielt. Danach las ich den Zettel am Schwarzen Brett. Es war eine Anzeige für ein Fahrrad. Ich las sie dreimal.

Herr Vogler kam an mir vorbei, sah kurz zu mir und sagte nichts. Ich sagte auch nichts. Wir haben es beide verstanden. Seitdem grüßen wir uns nicht mehr. Es funktioniert hervorragend.

\vfill
\begin{glosstbl}
  \textbf{das Stockwerk} & floor, storey & \textbf{der Schritt} & step, footstep \\
  \textbf{grüßen} & to greet, to say hello & \textbf{die Stufe} & stair, step \\
  \textbf{absichtlich} & deliberately, on purpose & \textbf{übersehen} & to overlook, to fail to notice \\
  \textbf{der Absatz} & landing (in a stairwell) & \textbf{das Schwarze Brett} & noticeboard \\
  \textbf{die Anzeige} & advertisement, notice & \textbf{hervorragend} & excellent, excellently \\
\end{glosstbl}

\gsec{Tagesordnungspunkt vier}
\levelbadge{B2}

Protokoll der Eigentümerversammlung, Tagesordnungspunkt vier: die Fußmatte vor Wohnung 3. Der Punkt war bereits im März vertagt worden. Die Sitzung dauerte an dieser Stelle einundfünfzig Minuten.

Frau Bittner erklärte, die Matte sei zu groß und rage in den Bereich hinein, der allen gehöre. Herr Kessler, dem die Matte gehört, entgegnete, sie sei bei seinem Einzug bereits vorhanden gewesen und werde daher von der Hausordnung nicht erfasst. Daraufhin wurde ein Zollstock geholt. Die Matte wurde in Anwesenheit aller gemessen. Das Ergebnis steht fest: sechzig mal vierzig Zentimeter.

Man einigte sich schließlich darauf, dass die Matte bleiben darf, solange sie gerade liegt. Wer sie künftig verschiebt, ist verpflichtet, sie zurückzuschieben. Nachtrag: Herr Kessler ist im Mai ausgezogen. Die Matte liegt noch da. Tagesordnungspunkt vier bleibt offen.

\vfill
\begin{glosstbl}
  \textbf{der Tagesordnungspunkt} & item on the agenda & \textbf{die Fußmatte} & doormat \\
  \textbf{vertagen} & to adjourn, to postpone & \textbf{hineinragen} & to jut into, to stick out into \\
  \textbf{entgegnen} & to counter, to reply & \textbf{der Einzug} & moving in \\
  \textbf{die Hausordnung} & house rules & \textbf{erfassen} & to cover, to take in \\
  \textbf{der Zollstock} & folding ruler & \textbf{sich einigen auf} & to reach agreement on \\
\end{glosstbl}

\gpart{XIX}{Geschichte}

\gsec{Der Zettel}
\levelbadge{A2}

Am 9. November 1989 gibt Günter Schabowski in Ost-Berlin eine Pressekonferenz. Er ist Mitglied der Regierung, aber er hat an der wichtigen Sitzung nicht teilgenommen. Kurz vorher drückt ihm ein Kollege einen Zettel in die Hand. Schabowski liest ihn nicht.

Am Ende liest er den Zettel laut vor: Die Bürger der DDR dürfen ab jetzt ins Ausland reisen. Ein Journalist fragt: \glqq{}Wann tritt das in Kraft?\grqq{} Schabowski sucht in seinen Papieren, findet nichts und sagt dann: \glqq{}Das tritt nach meiner Kenntnis ... sofort, unverzüglich.\grqq{}

Das stimmt nicht. Die neue Regel soll erst am nächsten Morgen gelten. Aber Millionen Menschen sehen die Konferenz im Fernsehen, und sie hören nur ein Wort: sofort. Kurz vor Mitternacht stehen Tausende an der Grenze. Ein Offizier ruft seine Vorgesetzten an, aber niemand kann ihm sagen, was er tun soll. Also öffnet er die Schranke. Die Mauer fällt, weil ein Mann seinen Zettel nicht gelesen hat.

\vfill
\begin{glosstbl}
  \textbf{der Zettel} & note, slip of paper & \textbf{teilnehmen an} & to take part in \\
  \textbf{in Kraft treten} & to come into force & \textbf{unverzüglich} & immediately, without delay \\
  \textbf{gelten} & to be valid, to apply & \textbf{der Vorgesetzte} & superior, boss \\
  \textbf{die Schranke} & barrier & \textbf{aus Versehen} & by accident, by mistake \\
\end{glosstbl}

\gsec{Die Uniform}
\levelbadge{B1}

Wilhelm Voigt war Schuster und hatte den größten Teil seines Lebens im Gefängnis verbracht. Als er 1906 freikam, brauchte er einen Pass. Ohne Pass bekam er keine Arbeit, und ohne Arbeit durfte er nirgends wohnen. Die Behörden blieben höflich und unerbittlich. Also kaufte er sich in Trödelläden Stück für Stück eine gebrauchte Offiziersuniform.

Am 16. Oktober 1906 stellte er sich in Berlin einer Gruppe Soldaten in den Weg und befahl ihnen, ihm zu folgen. Niemand verlangte Papiere, denn er trug ja eine Uniform. Mit der Bahn fuhr die kleine Truppe nach Köpenick. Dort besetzte Voigt das Rathaus, ließ den Bürgermeister verhaften und beschlagnahmte die Stadtkasse.

Dann verlangte er, was er die ganze Zeit gewollt hatte: einen Pass. Doch ausgerechnet das Rathaus von Köpenick hatte kein Passamt. Voigt nahm das Geld und verschwand. Die Zeitungen nannten ihn bald den Hauptmann von Köpenick, und halb Europa lachte. Zehn Tage später wurde er gefasst und zu vier Jahren Haft verurteilt. 1908 begnadigte ihn der Kaiser persönlich – also genau der Mann, dessen Uniform dem Schuster alle Türen geöffnet hatte.

\vfill
\begin{glosstbl}
  \textbf{der Schuster} & cobbler, shoemaker & \textbf{freikommen} & to be released, to get out \\
  \textbf{unerbittlich} & unrelenting, implacable & \textbf{der Trödelladen} & second-hand shop, junk shop \\
  \textbf{beschlagnahmen} & to confiscate, to seize & \textbf{ausgerechnet} & of all places, of all things \\
  \textbf{das Passamt} & passport office & \textbf{fassen} & to catch, to apprehend \\
  \textbf{verurteilen} & to sentence, to convict & \textbf{begnadigen} & to pardon \\
\end{glosstbl}

\gsec{Die Warnung}
\levelbadge{B2}

Im Jahr 9 nach Christus verwaltete Publius Quinctilius Varus die junge römische Provinz in Germanien. Sein wichtigster Verbündeter war Arminius, ein cheruskischer Fürst, der in Rom ausgebildet worden war und das römische Bürgerrecht besaß. Varus vertraute ihm vollkommen.

Bei einem Gastmahl nahm ihn ein anderer Cherusker beiseite. Segestes erklärte, Arminius plane einen Aufstand; Varus solle ihn und die übrigen Fürsten sofort festnehmen lassen – und ihn, Segestes, gleich mit, damit kein Unschuldiger falsch beurteilt werde. Varus hielt das für Neid unter Nachbarn und aß weiter.

Kurz darauf führte Arminius die Legionen in ein unwegsames Waldgebiet, aus dem es kein Entkommen gab. Drei Legionen wurden über mehrere Tage aufgerieben; Varus stürzte sich in sein Schwert. In Rom soll Augustus monatelang mit dem Kopf gegen die Türen geschlagen und dabei immer wieder gerufen haben: \glqq{}Varus, gib mir meine Legionen zurück!\grqq{} Die Nummern XVII, XVIII und XIX hat Rom nie wieder vergeben.

\vfill
\begin{glosstbl}
  \textbf{verwalten} & to govern, to administer & \textbf{der Verbündete} & ally \\
  \textbf{der Fürst} & prince, chieftain & \textbf{das Gastmahl} & banquet, feast \\
  \textbf{der Aufstand} & uprising, revolt & \textbf{festnehmen} & to arrest \\
  \textbf{der Neid} & envy & \textbf{unwegsam} & trackless, hard to cross \\
  \textbf{aufreiben} & to wear down, to wipe out & \textbf{vergeben} & to assign, to allocate \\
\end{glosstbl}

\gpart{XX}{Philosophie}

\gsec{Was übrig bleibt}
\levelbadge{B1}

Im Winter 1619 saß Descartes allein in einem warmen Zimmer und beschloss, alles zu bezweifeln, was er zu wissen glaubte. Seine Augen hatten ihn schon oft getäuscht, also traute er ihnen nicht mehr. Im Traum hatte er ebenso fest geglaubt, dass er wach war. Warum sollte er es jetzt besser wissen? Vielleicht, so dachte er, gab es sogar einen bösen Geist, der ihm die ganze Welt nur vorspielte.

Am Ende blieb ihm nichts. Der Himmel, sein Körper, sogar die Zahlen — alles konnte falsch sein. Doch dann fiel ihm etwas auf: Wenn ihn jemand täuscht, muss es jemanden geben, der getäuscht wird. Solange er zweifelte, war er da. \glqq{}Ich denke, also bin ich\grqq{} ist deshalb kein gewöhnlicher Beweis, sondern ein Satz, der wahr wird, sobald jemand ihn denkt.

Ob damit wirklich ein \glqq{}Ich\grqq{} bewiesen ist, haben spätere Philosophen bestritten; vielleicht zeigt der Gedanke nur, dass gerade gedacht wird. Versuchen Sie es trotzdem einmal ernsthaft: Denken Sie den Satz \glqq{}Ich existiere nicht\grqq{}. Es geht nicht — und Sie merken sofort, warum.

\vfill
\begin{glosstbl}
  \textbf{bezweifeln} & to doubt, to call into question & \textbf{trauen} & to trust (+ dative) \\
  \textbf{täuschen} & to deceive & \textbf{jemandem etwas vorspielen} & to fake something for someone \\
  \textbf{auffallen} & to strike someone, to be noticed & \textbf{der Beweis} & proof \\
  \textbf{gewöhnlich} & ordinary & \textbf{bestreiten} & to dispute, to deny \\
  \textbf{ernsthaft} & seriously & \textbf{sofort} & immediately \\
\end{glosstbl}

\gsec{Ein kleiner Kredit}
\levelbadge{B2}

Sie brauchen dringend Geld und wissen genau, dass Sie es nie zurückzahlen können. Ein Bekannter leiht es Ihnen nur, wenn Sie ihm versprechen, alles nächsten Monat zurückzugeben. Also versprechen Sie es. Kant fragt nun nicht, ob dadurch jemand geschädigt wird, und auch nicht, ob Sie ein schlechtes Gewissen haben. Er stellt eine ganz andere Frage.

Stellen Sie sich vor, Ihre Regel gelte für alle: Wer Geld braucht, verspricht, was er nicht halten kann. Wäre das der Fall, würde niemand mehr einem Versprechen glauben. Dann aber ließe sich mit einem Versprechen überhaupt kein Geld mehr beschaffen — die Handlung, die Sie planen, wäre in einer solchen Welt schlicht unmöglich. Das ist Kants Pointe: Die Lüge scheitert nicht an ihren Folgen, sondern an sich selbst. Damit Ihr Trick funktioniert, müssen Sie eine Ausnahme von genau der Regel machen, von der Sie leben.

Ob dieser Test wirklich jede moralische Frage entscheidet, wird bis heute bestritten. Der Kern bleibt aber bestehen: Manche Handlungen sind nicht deshalb falsch, weil sie schaden, sondern weil sie nur funktionieren, solange fast niemand sie begeht.

\vfill
\begin{glosstbl}
  \textbf{zurückzahlen} & to pay back & \textbf{das Gewissen} & conscience \\
  \textbf{schädigen} & to harm, to do damage to & \textbf{gelten für} & to apply to, to hold for \\
  \textbf{beschaffen} & to obtain, to procure & \textbf{die Pointe} & the point, the punchline \\
  \textbf{scheitern an} & to fail because of & \textbf{die Ausnahme} & exception \\
  \textbf{bestreiten} & to dispute, to contest & \textbf{begehen} & to commit (an act) \\
\end{glosstbl}

\gpart{XXI}{Wissenschaft}

\gsec{Die Uhr aus Licht}
\levelbadge{B2}

Stell dir eine Uhr vor, die nur aus zwei Spiegeln besteht. Zwischen ihnen springt ein Lichtblitz auf und ab; jede Rückkehr ist ein Tick. Steht die Uhr still, läuft das Licht senkrecht — den kürzesten Weg.

Nun stellst du dieselbe Uhr in einen fahrenden Zug. Vom Bahnsteig aus gesehen springt der Blitz nicht mehr senkrecht, sondern schräg, denn der obere Spiegel ist inzwischen weitergefahren. Der Weg ist also länger. Und hier kommt die eine Tatsache, die gemessen und nicht erfunden wurde: Licht ist für jeden Beobachter gleich schnell. Ein längerer Weg bei gleicher Geschwindigkeit bedeutet mehr Zeit pro Tick. Die Uhr im Zug geht langsamer.

Für die Reisende im Zug ist nichts seltsam: ihr Blitz springt senkrecht, ihre Uhr geht normal, und es ist der Bahnsteig, dessen Uhren nachgehen. Beide haben recht, solange der Zug gleichmäßig fährt. Und es ist keine Eigenart der Spiegel — Herzschlag, Zerfall, Gedanke: alles verlangsamt sich mit, weil es die Zeit selbst ist, die gedehnt wird.

\vfill
\begin{glosstbl}
  \textbf{der Spiegel} & mirror & \textbf{der Lichtblitz} & flash of light \\
  \textbf{senkrecht} & vertical, straight up & \textbf{schräg} & slanted, at an angle \\
  \textbf{der Zerfall} & decay & \textbf{nachgehen} & to run slow (of a clock) \\
  \textbf{die Eigenart} & quirk, peculiarity & \textbf{dehnen} & to stretch \\
  \textbf{gleichmäßig} & steady, even &  &  \\
\end{glosstbl}

\gsec{Der zweite Pass}
\levelbadge{B1}

Zwei Täler, dazwischen ein Gebirge. Im linken Tal stehen die Ausgangsstoffe, im rechten liegt das Produkt — tiefer, also günstiger. Trotzdem passiert nichts, denn zuerst muss jedes Molekül über den Kamm. Nur wenige haben dafür genug Schwung. Die Reaktion läuft, aber quälend langsam.

Ein Katalysator schiebt nun niemanden über den Berg. Er baut auch das rechte Tal nicht tiefer. Er öffnet einen zweiten, niedrigeren Pass. Plötzlich schaffen es viel mehr Moleküle hinüber, und was Jahre gedauert hätte, ist in Sekunden vorbei.

Zwei Dinge folgen daraus. Erstens kommt der Katalysator am Ende unverändert wieder heraus — er macht mit, aber er wird nicht verbraucht, deshalb genügen winzige Mengen. Zweitens ändert er nur das Tempo, nie das Ziel. Liegt das rechte Tal höher, bleiben die meisten Moleküle links liegen — daran ändert der schönste Pass nichts. Ein Katalysator macht das Mögliche schnell — er macht das Unmögliche nicht möglich.

\vfill
\begin{glosstbl}
  \textbf{das Tal} & valley & \textbf{das Gebirge} & mountain range \\
  \textbf{der Ausgangsstoff} & starting material, reactant & \textbf{der Kamm} & ridge, crest \\
  \textbf{der Schwung} & momentum, energy & \textbf{quälend} & agonisingly \\
  \textbf{es hinüberschaffen} & to make it across & \textbf{verbrauchen} & to use up \\
  \textbf{winzig} & tiny & \textbf{genügen} & to be enough \\
\end{glosstbl}

